\small
\begin{longtable}{@{}>{\raggedright\arraybackslash}p{.05\textwidth}>{\raggedright\arraybackslash}p{.40\textwidth}>{\raggedright\arraybackslash}p{.49\textwidth}@{}}
\toprule No. & Question in the fixed edition & Disposition and qualification\\\midrule\endhead
1 & Classify self-dual toroidal polyhedra. & Unsettled here. The spherical classification is attributed background; regular or chiral toroidal maps are only a subclass.\\[5pt]
2 & Decide whether every equivelar self-dual polyhedron has a polarity. & Unsettled here. The source gives the affirmative regular/chiral subclass and a negative unrestricted self-duality example; neither answers the equivelar question.\\[5pt]
3 & Classify fully transitive polyhedra on surfaces of small genus, with both orientability choices. & Unsettled here. The toroidal classification cited in the source includes a further family whose middle adjacency preserves each flag orbit.\\[5pt]
4 & Determine the least orientable genus at least two admitting a nonregular fully transitive polyhedron. & Unsettled here. The source excludes chiral maps in genera two through six; fully transitive is a larger class.\\[5pt]
5 & Characterize groups of fully transitive polyhedra outside the regular and chiral classes. & Unsettled here. The subsequent self-dual restriction is an additional suggested subproblem.\\[5pt]
6 & Classify edge-transitive polyhedra on orientable and nonorientable surfaces of small genus. & Unsettled here. The cited partial results do not give the requested classification.\\[5pt]
7 & Determine what all specified hole and zigzag lengths determine, and when the associated universal regular quotient is finite. & Unsettled here. The first question covers regular or chiral polyhedra of type \{p,q\}, with 1<=j<=floor(q/2). The universal regular quotient in the subsequent question is distinct. Word-order formulas alone give neither uniqueness nor a universal finiteness criterion.\\[5pt]
8 & Decide existence with prescribed regular or chiral facets and vertex figures. & Already negative in the source. This is already negative in general in the source: the cited toroidal (3,0) facet and hemicube vertex figure have no regular amalgamation.\\[5pt]
9 & Characterize the compatible facet and vertex-figure pairs for which amalgamation exists. & Unsettled here. Matching their common sections is necessary but insufficient.\\[5pt]
10 & Determine universal groups and their finiteness for regular facet and vertex-figure pairs admitting amalgamation. & Unsettled here. The source existence theorem for the universal object is retained separately from group classification.\\[5pt]
11 & Determine universal groups and finiteness for oriented chiral/regular amalgamations with at least one chiral member. & Unsettled here. Choose enantiomorphic forms, with at least one chiral constituent and every regular constituent directly regular. Simultaneous reversal identifies four oriented choices in pairs; if one constituent is regular the remaining choices also agree. Existence remains a prerequisite.\\[5pt]
12 & Determine existence and finiteness with an arbitrary finite regular facet and a convex regular vertex figure. & Unsettled here. The source records the cross-polytope vertex-figure case: finiteness is equivalent to neighborliness of the facet.\\[5pt]
13 & Determine existence and finiteness with finite chiral facets and convex regular vertex figures. & Unsettled here. The regular cross-polytope result does not automatically extend to chiral facets.\\[5pt]
14 & Establish existence of finite chiral polytopes in ranks at least six. & Unsettled here. Affirmative answer conjectured in the 2006 source. This inventory makes no present-day open-status assertion.\\[5pt]
15 & Decide whether a convex four-polytope can have every facet combinatorially icosahedral. & Unsettled here. An abstract polytope or incidence geometry with such facets is not sufficient.\\[5pt]
16 & Classify universal regular amalgamations of square toroidal maps with odd distinct parameters s,u at least three. & Unsettled here. Two claims are embedded: existence for every such pair and finiteness only for (3,5) and (5,3).\\[5pt]
17 & Classify universal regular amalgamations of triangular and hexagonal toroidal maps in the stated regular parameter families. & Unsettled here. Retain the exact source parameter domains: s>=2,t=0 or s=t>=1, and the corresponding u,v alternatives.\\[5pt]
18 & Complete the classification of the three locally toroidal universal regular rank-six families, up to duality. & Unsettled here. The completeness of each finite list and the explicit one-relation infinite-family conjecture are separate clauses.\\[5pt]
19 & Classify the universal locally toroidal chiral rank-four polytopes with the specified regular or chiral constituents and enantiomorphic forms. & Unsettled here. The preceding discussion and subsequent group calculation focus on the finite classification; retain that goal separately from the literal universal-object wording. At least one constituent is a chiral toroidal map. The square-torus/cube subfamily has its own four-pair finiteness conjecture.\\[5pt]
20 & Classify finite universal regular or chiral rank-four polytopes with prescribed small facet and vertex-figure genera and orientabilities. & Unsettled here. Duality exchanges the two mixed orientability cases; the problem is not restricted to locally toroidal polytopes.\\[5pt]
21 & Develop the realization theory of chiral polytopes. & Unsettled here. A theory of equivariant Gram matrices alone does not classify faithful realizations or determine their full geometric symmetry.\\[5pt]
22 & Describe realization cones of regular toroids in every rank. & Unsettled here. The cited rank-three cones are established background, not the requested all-rank description.\\[5pt]
23 & Decide whether a finite string C-group can have a complex irreducible character with Frobenius-Schur indicator minus one. & Unsettled here. The representation must be quaternionic; a general real representation or a non-C-group example is insufficient.\\[5pt]
24 & Determine intersection-property validity from character tables or other readily computed group-algebra invariants for a marked string-generated group. & Unsettled here. The marking is part of the input. An explicit subgroup enumeration is not by itself an answer to the character-data question.\\[5pt]
25 & Enumerate discrete fully transitive polyhedra in Euclidean three-space, both finite and infinite. & Unsettled here. Regular and chiral classifications supply subclasses only.\\[5pt]
26 & Enumerate discrete uniform polyhedra in Euclidean three-space, both finite and infinite. & Unsettled here. The finite planar-faced list does not include all skew-faced or infinite cases.\\[5pt]
27 & Characterize when the source construction from a finite bipartite trivalent two- or three-arc-transitive graph yields a polytope. & Unsettled here. The condition must concern the actual reconstructed poset; graph transitivity alone is not asserted sufficient.\\[5pt]
28 & Decide whether all Petrie schemes are acoptic for each of four specified classes of products, uniform convex polytopes and regular star-polytopes. & Unsettled here. The four classes are: products of cubes and crosspolytopes; products of cubes or crosspolytopes and 4k-gons (k>=1); uniform convex polytopes; regular star-polytopes. Affirmative answers are conjectured for classes 1,2,4; class 3 remains a separate question.\\[5pt]
29 & Characterize universal regular amalgamations for which every Petrie scheme is acoptic. & Unsettled here. The stated local nonintersection condition on constituents is only necessary.\\[5pt]
30 & Characterize full automorphism groups of regular or chiral polytopes of order a power of two or a power of two times an odd prime. & Restricted acquired results only; general question unsettled here. Existing project results cover specified ordered regular markings, bounded class/excess and exact small excesses; the unmarked, chiral and mixed-prime problem remains outside them.\\[5pt]
31 & Find polytopes with sporadic simple full automorphism groups, explicitly including the Monster. & Affirmative existence by the acquired 2016 theorem. The accepted positive theorem of the 2016 simple-group manuscript gives a chiral polyhedron for each of the26 sporadic groups. This is an acquired earlier answer, not a new conjecture proof.\\[5pt]
32 & Find regular, chiral, or other polytopes with full automorphism group A\_n; in particular determine the possible degrees for each rank in the regular or chiral classes. & Regular ranks acquired; chiral rank-three existence only; other parts unsettled here. For n>=3, acquired L038/L044 and their complete small censuses settle the regular ranks: A5:\{3\}; A9:\{3,4\}; A10:\{3,4,5\}; A11:\{3,6\}; none at n=3,4,6,7,8; and exactly 3<=r<=floor((n-1)/2) for n>=12. Acquired L045 gives chiral rank-three existence for n>=8. This does not classify all chiral ranks or all other polytopes. These are attributed earlier results, not newly solved questions.\\[5pt]
33 & Classify finite neighborly regular or chiral polytopes. & Unsettled here. The necessary two-transitive vertex action is not asserted sufficient for reconstruction.\\[5pt]
34 & Decide whether semisparsity follows from semisparsity of both end intersections in every conjugate and the displayed product-intersection equality. & Unsettled here. Retain every conjugate quantifier; necessity and the general three-condition definition do not prove this sufficiency assertion.\\[5pt]
\bottomrule
\end{longtable}
\normalsize
